\begin{textsample}{POS Dim 2 – vem\_ed\_19 – Score 31.00 – vem\_ed\_19/t088.txt}  \label{ex:f2_pos_002}
Que conselho você daria para professores interessados em \textbf{manter} Intercâmbios Virtuais por um longo tempo?
Tentamos fazer os professores \textbf{enxergarem} oportunidades que os projetos COIL oferecem para além da \textbf{duração} das colaborações.
Se é uma parceria nova, encorajamos a oferecer o mesmo projeto por vários \textbf{semestres}.
O planejamento da primeira edição leva mais tempo, mas da segunda \textbf{oferta} em diante, são mudanças menores.
Utilizar a mesma estrutura de projeto com outros grupos de estudantes, e com parceiros em diversos países, traz \textbf{múltiplas} perspectivas.
Outro aspecto que \textbf{vejo} crescer é a pesquisa: professores trabalhando juntos \textbf{investigam} a metodologia de Intercâmbios Virtuais em seu \textbf{campo} de conhecimento e em suas turmas.
E publicam artigos em \textbf{periódicos} acadêmicos sobre essa \textbf{inovação}.
Se têm parceiros em vários países, podem realizar pesquisas \textbf{multidisciplinares} em diversas regiões do mundo.
Há muita \textbf{demanda} para pesquisa sobre Intercâmbios Virtuais: sobre o desenvolvimento de habilidades dos alunos, sobre a aprendizagem sob as perspectivas comparadas das turmas e como isso muda de acordo com o país — os resultados são \textbf{parecidos} ou diferentes?
Há tantas oportunidades para produzir artigos, \textbf{periódicos} ou até mesmo \textbf{capítulos} de \textbf{livros}.
Qual sua visão para os Intercâmbios Virtuais nos próximos cinco anos?
Talvez por \textbf{causa} de meu papel, \textbf{vejo} todos os benefícios que os Intercâmbios Virtuais podem trazer. \textbf{Acho} que COIL veio para \textbf{ficar} e vai crescer à medida que mais instituições se \textbf{juntarem} a essa comunidade.
A pandemia ajudou de certa forma as instituições que buscam oferecer experiências de internacionalização para seus estudantes.
Parte de minha visão e \textbf{esperança} é que haja mais apoio \textbf{estável} para iniciativas de Intercâmbio Virtual nas instituições.
Muitas começam sem equipe, apenas com uma pessoa que vê COIL como uma grande oportunidade para seus professores, estudantes e para a instituição como um todo.
Depois de um tempo, conseguem mais recursos.
É importante reconhecer o papel das estruturas de \textbf{suporte} nas instituições - tecnológico, de \textbf{design} de projetos, de parcerias, linguístico. \textbf{Acho} que o desenvolvimento da tecnologia de Inteligência Artificial para auxiliar nas \textbf{traduções} simultâneas \textbf{deve} \textbf{facilitar} o trabalho colaborativo, especialmente o de pequenos grupos de estudantes.
Também \textbf{preciso} mencionar as redes de coordenadores de Intercâmbios Virtuais nas diferentes instituições.
É importante aprender uns com os outros, porque há um conhecimento \textbf{incrível} nas diferentes regiões do mundo, e \textbf{manter} vivas essas redes.
É uma comunidade \textbf{maravilhosa}, que pode se \textbf{beneficiar} de recursos e ferramentas que ajudem a todas as instituições.

% matched lemmas: achar, beneficiar, campo, capítulo, causa, demanda, design, dever, duração, enxergar, esperança, estável, facilitar, ficar, incrível, inovação, investigar, juntar, livro, manter, maravilhoso, multidisciplinar, múltiplo, oferta, parecir, periódico, preciso, semestr, suporte, tradução, vejo
\end{textsample}
